% Options for packages loaded elsewhere
\PassOptionsToPackage{unicode}{hyperref}
\PassOptionsToPackage{hyphens}{url}
\PassOptionsToPackage{dvipsnames,svgnames,x11names}{xcolor}
\documentclass[
  11pt,
]{article}
\usepackage{xcolor}
\usepackage[margin=2.5cm]{geometry}
\usepackage{amsmath,amssymb}
\setcounter{secnumdepth}{5}
\usepackage{iftex}
\ifPDFTeX
  \usepackage[T1]{fontenc}
  \usepackage[utf8]{inputenc}
  \usepackage{textcomp} % provide euro and other symbols
\else % if luatex or xetex
  \usepackage{unicode-math} % this also loads fontspec
  \defaultfontfeatures{Scale=MatchLowercase}
  \defaultfontfeatures[\rmfamily]{Ligatures=TeX,Scale=1}
\fi
\usepackage{lmodern}
\ifPDFTeX\else
  % xetex/luatex font selection
\fi
% Use upquote if available, for straight quotes in verbatim environments
\IfFileExists{upquote.sty}{\usepackage{upquote}}{}
\IfFileExists{microtype.sty}{% use microtype if available
  \usepackage[]{microtype}
  \UseMicrotypeSet[protrusion]{basicmath} % disable protrusion for tt fonts
}{}
\makeatletter
\@ifundefined{KOMAClassName}{% if non-KOMA class
  \IfFileExists{parskip.sty}{%
    \usepackage{parskip}
  }{% else
    \setlength{\parindent}{0pt}
    \setlength{\parskip}{6pt plus 2pt minus 1pt}}
}{% if KOMA class
  \KOMAoptions{parskip=half}}
\makeatother
\usepackage{graphicx}
\makeatletter
\newsavebox\pandoc@box
\newcommand*\pandocbounded[1]{% scales image to fit in text height/width
  \sbox\pandoc@box{#1}%
  \Gscale@div\@tempa{\textheight}{\dimexpr\ht\pandoc@box+\dp\pandoc@box\relax}%
  \Gscale@div\@tempb{\linewidth}{\wd\pandoc@box}%
  \ifdim\@tempb\p@<\@tempa\p@\let\@tempa\@tempb\fi% select the smaller of both
  \ifdim\@tempa\p@<\p@\scalebox{\@tempa}{\usebox\pandoc@box}%
  \else\usebox{\pandoc@box}%
  \fi%
}
% Set default figure placement to htbp
\def\fps@figure{htbp}
\makeatother
% definitions for citeproc citations
\NewDocumentCommand\citeproctext{}{}
\NewDocumentCommand\citeproc{mm}{%
  \begingroup\def\citeproctext{#2}\cite{#1}\endgroup}
\makeatletter
 % allow citations to break across lines
 \let\@cite@ofmt\@firstofone
 % avoid brackets around text for \cite:
 \def\@biblabel#1{}
 \def\@cite#1#2{{#1\if@tempswa , #2\fi}}
\makeatother
\newlength{\cslhangindent}
\setlength{\cslhangindent}{1.5em}
\newlength{\csllabelwidth}
\setlength{\csllabelwidth}{3em}
\newenvironment{CSLReferences}[2] % #1 hanging-indent, #2 entry-spacing
 {\begin{list}{}{%
  \setlength{\itemindent}{0pt}
  \setlength{\leftmargin}{0pt}
  \setlength{\parsep}{0pt}
  % turn on hanging indent if param 1 is 1
  \ifodd #1
   \setlength{\leftmargin}{\cslhangindent}
   \setlength{\itemindent}{-1\cslhangindent}
  \fi
  % set entry spacing
  \setlength{\itemsep}{#2\baselineskip}}}
 {\end{list}}
\usepackage{calc}
\newcommand{\CSLBlock}[1]{\hfill\break\parbox[t]{\linewidth}{\strut\ignorespaces#1\strut}}
\newcommand{\CSLLeftMargin}[1]{\parbox[t]{\csllabelwidth}{\strut#1\strut}}
\newcommand{\CSLRightInline}[1]{\parbox[t]{\linewidth - \csllabelwidth}{\strut#1\strut}}
\newcommand{\CSLIndent}[1]{\hspace{\cslhangindent}#1}
\setlength{\emergencystretch}{3em} % prevent overfull lines
\providecommand{\tightlist}{%
  \setlength{\itemsep}{0pt}\setlength{\parskip}{0pt}}
\usepackage{bookmark}
\IfFileExists{xurl.sty}{\usepackage{xurl}}{} % add URL line breaks if available
\urlstyle{same}
% fallback for those not using the hyperref driver hyperxmp:
\makeatletter
\@ifundefined{xmpquote}{\newcommand{\xmpquote}[1]{#1}}{}
\makeatother
\hypersetup{
  pdftitle={SCORE: Synthesising Complete{,} Owned{,} Responsive Environments for Simulation},
  pdfauthor={J.P. Kardolus, Babon Innovations B.V., Utrecht, The Netherlands},
  colorlinks=true,
  linkcolor={Maroon},
  filecolor={Maroon},
  citecolor={Blue},
  urlcolor={Blue},
  pdfcreator={LaTeX via pandoc}}

\title{SCORE: Synthesising Complete, Owned, Responsive Environments for
Simulation}
\author{J.P. Kardolus, Babon Innovations B.V., Utrecht, The Netherlands}
\date{Working draft, 18 November 2026}

\begin{document}
\maketitle
\begin{abstract}
A single locked concept picture becomes an owned, editable OpenUSD
environment that a physics engine can step, populated with people solved
down to the muscle. World generation models now render convincing
pictures, videos and Gaussian splats of places, yet none delivers
separate, editable objects with authored mass, colliders and a
consistent finish that have been checked in simulation. SCORE is an open
framework that composes current commercial and open-source models around
the open OpenUSD standard and the open-source MuJoCo engine. From a
creator's concept it segments the picture into objects, generates each
missing one, finishes every object from a shared library, checks the
assembled scene for physical faults and populates it with people whose
MyoFullBody musculoskeletal bodies are solved by MOTOR. Four places
built this way pass every check without a fault and are written as
stages their creator owns and can hand to a simulator as written; the
people remain the open layer, since their sim bodies do not yet meet the
residual limits over whole clips. Code:
https://github.com/Babon-Innovations-b-v/score (MIT).
\end{abstract}

\makeatletter\def\fps@figure{tp}\makeatother\renewcommand{\topfraction}{0.9}\renewcommand{\textfraction}{0.1}

\section{Related work}\label{related-work}

OpenUSD began at Pixar in 2012 as the scene description for assembling
its films {[}1{]} and was released as open source in 2016 {[}2{]}. A USD
stage is composed from layers, references and variants, so that each
tool edits its own layer. In 2023 Pixar, Adobe, Apple, Autodesk and
NVIDIA founded the Alliance for OpenUSD to standardise it {[}3{]}; its
Core Specification 1.0, ratified in December 2025, covers composition
and data representation and leaves physics to later domain
specifications {[}4{]}. The UsdPhysics schema already describes rigid
bodies, colliders, mass and joints {[}5{]}, and Isaac Lab simulates such
stages for robot learning {[}6{]}, but every object's mass, collider and
joints must still be authored.

MuJoCo {[}7{]}, built on generalised coordinates and soft contact, has
been open source since 2022 {[}8{]} and is widely used in robotics and
musculoskeletal simulation {[}9{]}. MuJoCo Warp ports it to the GPU
{[}10{]}, {[}11{]}, and Newton, an open engine under the Linux
Foundation since September 2025, runs it as its main solver and reads
OpenUSD directly {[}12{]}. MuscleMimic assembles MyoFullBody, a
head-to-toe body with 416 muscles, from MyoSuite's models {[}13{]}. Each
still needs a scene whose physics has been authored.

World generation models do not supply such scenes. Genie 3 generates
interactive worlds as pixels from a closed model {[}14{]} and Cosmos
predicts video {[}15{]}; neither holds exportable geometry or physical
state. HunyuanWorld 1.0 lifts panoramas into layered meshes with baked
colour and no physical properties {[}16{]}, and its successor
reconstructs splats and meshes with coarse collision but no per-object
mass or joints {[}17{]}. Object generators such as TRELLIS.2 {[}18{]}
produce single assets with no scale, mass or relation to a scene, and
procedural and language-driven systems keep scenes explicit but draw on
a fixed grammar or asset set {[}19{]}, {[}20{]}, {[}21{]}. What is
missing is a framework that composes current commercial and open-source
components into one route from a concept to a scene a simulator can
step.

\section{Method}\label{method}

\begin{figure}
\centering
\includegraphics[width=1\linewidth,height=\textheight,keepaspectratio,alt={One
place, camp, through SCORE's layers. (a) The creator's locked concept.
(b) SAM 3 segments it into 80 instances from its probe nouns. (c) The
two objects of camp made by TRELLIS.2, each close-up beside its
generated mesh. (d) The assembled place, every surface finished from the
place's library palette (strip, 31 tokens). (e) CoACD convex
decomposition written with the model (43 hulls for the desk, 48 for the
door), with mass, friction and restitution from the material class. (f)
A walker's MyoFullBody sim body, solved by MOTOR, in the finished
place.}]{figures/layers.jpg}
\caption{One place, camp, through SCORE's layers. (a) The creator's
locked concept. (b) SAM 3 segments it into 80 instances from its probe
nouns. (c) The two objects of camp made by TRELLIS.2, each close-up
beside its generated mesh. (d) The assembled place, every surface
finished from the place's library palette (strip, 31 tokens). (e) CoACD
convex decomposition written with the model (43 hulls for the desk, 48
for the door), with mass, friction and restitution from the material
class. (f) A walker's MyoFullBody sim body, solved by MOTOR, in the
finished place.}
\end{figure}

SCORE writes every scene as an OpenUSD stage in metres, its physics in
UsdPhysics and its child scenes by reference, so that each layer is
edited without disturbing the others. The creator first decides physics,
setting and style in an environment concept: its bodies and orbits, each
plot's ground and gravity, site coordinates, a style line and reference
pictures. A set interview with Gemini 3.1 Pro fills it {[}22{]} and the
creator locks it; each site's sky is computed from these bodies.

Figure 1 follows one place, camp, through the layers. Nano Banana Pro
draws its concept through the Google Gemini API {[}23{]} from the
place's contents, view and style, for the creator to lock. SAM 3
{[}24{]} segments the locked concept into its inventory from ten generic
nouns, keeping up to 40 instances each and merging overlaps above 0.6. A
coding agent {[}25{]} lays library entries over it, renders the view in
Blender {[}26{]} and corrects until it reads as the concept. TRELLIS.2
{[}18{]} makes objects the library lacks from a close-up or the masked
concept crop, fitting each over its placeholder. The library finish
moves each of six colour groups 60\% of the way to the nearest token of
the place's one palette, so that generated and library objects share a
look, and flags a group far from every token, which the agent repaints
with a library surface {[}27{]}. CoACD {[}28{]} writes convex hulls with
each model for downstream simulators.

The assembled scene is checked for twelve kinds of fault, most of them
geometric, among them resting more than 3 cm off a support,
interpenetration, tipping and a walkway narrower than 0.9 m. In MuJoCo,
every loose object is released for 1 s at a 2 ms step, and one that
moves more than 5 cm or turns more than 8° is unstable. Only a
fault-free scene is ready.

People are routed over the checked floor, animated by Kimodo {[}29{]}
and solved by MOTOR {[}30{]} as MyoFullBody sim bodies sized to their
SOMA bodies {[}31{]}. MOTOR fits the 95 free degrees of freedom of
SCORE's variant (131 joints, 36 tied), runs inverse dynamics in Newton
under the place's gravity with ground reaction forces from skin contact,
and shares the joint moments over 416 muscles by static optimisation. A
body passes the residual check when its residual force stays within 2\%
of body weight and its moment within 1\% of weight times centre-of-mass
height {[}32{]}.

\section{Discussion}\label{discussion}

\makeatletter\def\fps@figure{!htb}\makeatother\renewcommand{\bottomfraction}{0.7}

\begin{figure}
\centering
\includegraphics[width=1\linewidth,height=\textheight,keepaspectratio,alt={(a)--(d)
Opening frames of the four released films: camp, campgrounds (cropped),
street and stairwell. (e)--(g) One walker in the street film, peeled
from its player model to its SOMA body to its MyoFullBody sim body
solved by MOTOR, with the vertical ground reaction force drawn at the
stance foot.}]{figures/results.jpg}
\caption{(a)--(d) Opening frames of the four released films: camp,
campgrounds (cropped), street and stairwell. (e)--(g) One walker in the
street film, peeled from its player model to its SOMA body to its
MyoFullBody sim body solved by MOTOR, with the vertical ground reaction
force drawn at the stance foot.}
\end{figure}

SCORE first derived each scene from its concept through model readings,
written plans and generated shells. Over all 26 places of a first
environment, reviewing vision-model agents scored it a mean of 2.2 out
of 10 against the concepts, with wrong shells in 18 places and objects
dropped or doubled in 22, and eight rounds of fixes did not improve it.
The layered method replaced it for four places: camp, campgrounds,
stairwell and street (Figure 2).

All four pass the place check without a fault, on 711, 1046, 61 and 890
objects. The builder's own critic scored their concept views 7, 7, 4 and
6 out of 10, the old method's reviewers 2, 4, 1 and 1; with different
judges and no human rating, the two sets are not directly comparable.
TRELLIS.2 was judged better than the earlier route on five of six camp
props and runs in 170 to 220 s per model on a local 16 GB card at no
rental cost, against about EUR 0.22 before. First film renders cost
about EUR 17 on rented cards, the published re-renders EUR 45 to 50.

The people are the weakest layer. In the first films of 17 November,
twelve of thirteen walkers stayed within 6.8 cm of their routes, but no
sim body passed the residual check over a whole clip: the best street
walker met the force limit (13.2 N against 15.2 N) but not the moment
limit (13.4 N m against 7.4 N m), and the stairwell woman peaked at
1,368 N against 15 N while sinking about 17 cm. The published re-films
still fail it at every walker's peaks. The scenes are not
photorealistic: surfaces are clean and pale, kits repeat and signs are
illegible. The settle runs under Earth's gravity even on Mars, and the
check judges rest, not function.

A revised foot contact in MOTOR lowered per-foot vertical ground
reaction force error in all four motions of 13 trials but raised the
pelvis residual moment by 1.5 to 6 N m {[}33{]}, so the residual limits
remain open work. Articulated UsdPhysics joints would make doors and
machines operable, and the settle should take each plot's own gravity.
The larger step is a whole environment judged by human raters. Coding
agents placed every object through the framework's commands, judged the
renders and wrote most of the code; the creator's locks and decisions
were his own.

\section*{References}\label{references}
\addcontentsline{toc}{section}{References}

\protect\phantomsection\label{refs}
\begin{CSLReferences}{0}{0}
\bibitem[\citeproctext]{ref-openusd}
\CSLLeftMargin{{[}1{]} }%
\CSLRightInline{Pixar Animation Studios and Alliance for OpenUSD,
{``{OpenUSD}: Universal scene description.''} 2026. Available:
\url{https://openusd.org/release/intro.html}}

\bibitem[\citeproctext]{ref-pixar2016usd}
\CSLLeftMargin{{[}2{]} }%
\CSLRightInline{Pixar Animation Studios, {``Pixar releases universal
scene description as open source.''} Press release, Emeryville,
California, Jul. 2016. Available:
\url{https://openusd.org/dev/press_opensource_release.html}}

\bibitem[\citeproctext]{ref-aousd2023}
\CSLLeftMargin{{[}3{]} }%
\CSLRightInline{Pixar, Adobe, Apple, Autodesk, and NVIDIA, {``{Pixar,
Adobe, Apple, Autodesk, and NVIDIA form Alliance for OpenUSD to drive
open standards for 3D content}.''} Press release, Aug. 2023. Available:
\url{https://nvidianews.nvidia.com/news/aousd-to-drive-open-standards-for-3d-content}}

\bibitem[\citeproctext]{ref-aousd2025corespec}
\CSLLeftMargin{{[}4{]} }%
\CSLRightInline{Alliance for OpenUSD, {``{Alliance for OpenUSD Announces
Core Specification 1.0}.''} Dec. 2025. Available:
\url{https://aousd.org/news/core-spec-announcement/}}

\bibitem[\citeproctext]{ref-usdphysics}
\CSLLeftMargin{{[}5{]} }%
\CSLRightInline{Pixar Animation Studios and NVIDIA, {``{UsdPhysics}:
{USD} physics schema.''} 2026. Available:
\url{https://openusd.org/dev/api/usd_physics_page_front.html}}

\bibitem[\citeproctext]{ref-mittal2025isaaclab}
\CSLLeftMargin{{[}6{]} }%
\CSLRightInline{{NVIDIA \emph{et al.}}, {``{Isaac Lab}: A
{GPU}-accelerated simulation framework for multi-modal robot
learning,''} \emph{arXiv preprint arXiv:2511.04831}, 2025.}

\bibitem[\citeproctext]{ref-todorov2012mujoco}
\CSLLeftMargin{{[}7{]} }%
\CSLRightInline{E. Todorov, T. Erez, and Y. Tassa, {``{MuJoCo}: A
physics engine for model-based control,''} in \emph{IEEE/RSJ
international conference on intelligent robots and systems (IROS)},
2012, pp. 5026--5033.}

\bibitem[\citeproctext]{ref-mujocorepo}
\CSLLeftMargin{{[}8{]} }%
\CSLRightInline{Google DeepMind, {``{MuJoCo}: Multi-joint dynamics with
contact.''} \url{https://github.com/google-deepmind/mujoco}, 2022.}

\bibitem[\citeproctext]{ref-caggiano2022myosuite}
\CSLLeftMargin{{[}9{]} }%
\CSLRightInline{V. Caggiano, H. Wang, G. Durandau, M. Sartori, and V.
Kumar, {``{MyoSuite}: A contact-rich simulation suite for
musculoskeletal motor control,''} \emph{arXiv preprint
arXiv:2205.13600}, 2022.}

\bibitem[\citeproctext]{ref-mujocowarp}
\CSLLeftMargin{{[}10{]} }%
\CSLRightInline{Google DeepMind and NVIDIA, {``{MuJoCo Warp}: A
{GPU}-optimized version of the {MuJoCo} physics simulator.''} 2026.
Available: \url{https://github.com/google-deepmind/mujoco_warp}}

\bibitem[\citeproctext]{ref-warp}
\CSLLeftMargin{{[}11{]} }%
\CSLRightInline{M. Macklin, {``{Warp}: A high-performance {Python}
framework for {GPU} simulation and graphics.''} NVIDIA, 2022. Available:
\url{https://github.com/NVIDIA/warp}}

\bibitem[\citeproctext]{ref-newton}
\CSLLeftMargin{{[}12{]} }%
\CSLRightInline{The Newton Developers, {``{Newton}: A {GPU}-accelerated
physics simulation engine.''} 2026. Available:
\url{https://github.com/newton-physics/newton}}

\bibitem[\citeproctext]{ref-li2026musclemimic}
\CSLLeftMargin{{[}13{]} }%
\CSLRightInline{C. Li \emph{et al.}, {``Towards embodied {AI} with
{MuscleMimic}: Unlocking full-body musculoskeletal motor learning at
scale,''} \emph{arXiv preprint arXiv:2603.25544}, 2026.}

\bibitem[\citeproctext]{ref-genie3}
\CSLLeftMargin{{[}14{]} }%
\CSLRightInline{J. Parker-Holder and S. Fruchter, {``{Genie 3: A new
frontier for world models}.''} Google DeepMind blog, Aug. 2025.
Available:
\url{https://deepmind.google/blog/genie-3-a-new-frontier-for-world-models/}}

\bibitem[\citeproctext]{ref-agarwal2025cosmos}
\CSLLeftMargin{{[}15{]} }%
\CSLRightInline{{NVIDIA \emph{et al.}}, {``{Cosmos} world foundation
model platform for physical {AI},''} \emph{arXiv preprint
arXiv:2501.03575}, 2025.}

\bibitem[\citeproctext]{ref-hunyuanworld2025}
\CSLLeftMargin{{[}16{]} }%
\CSLRightInline{HunyuanWorld Team, {``{HunyuanWorld 1.0}: Generating
immersive, explorable, and interactive {3D} worlds from words or
pixels,''} \emph{arXiv preprint arXiv:2507.21809}, 2025.}

\bibitem[\citeproctext]{ref-hyworld2}
\CSLLeftMargin{{[}17{]} }%
\CSLRightInline{Tencent Hunyuan, {``{HY-World} 2.0.''} GitHub
repository, 2026. Available:
\url{https://github.com/Tencent-Hunyuan/HY-World-2.0}}

\bibitem[\citeproctext]{ref-xiang2025trellis2}
\CSLLeftMargin{{[}18{]} }%
\CSLRightInline{J. Xiang \emph{et al.}, {``Native and compact structured
latents for {3D} generation,''} \emph{arXiv preprint arXiv:2512.14692},
2025.}

\bibitem[\citeproctext]{ref-raistrick2023infinigen}
\CSLLeftMargin{{[}19{]} }%
\CSLRightInline{A. Raistrick \emph{et al.}, {``Infinite photorealistic
worlds using procedural generation,''} in \emph{IEEE/CVF conference on
computer vision and pattern recognition (CVPR)}, 2023.}

\bibitem[\citeproctext]{ref-raistrick2024indoors}
\CSLLeftMargin{{[}20{]} }%
\CSLRightInline{A. Raistrick \emph{et al.}, {``{Infinigen Indoors}:
Photorealistic indoor scenes using procedural generation,''} in
\emph{Proceedings of the IEEE/CVF conference on computer vision and
pattern recognition (CVPR)}, 2024.}

\bibitem[\citeproctext]{ref-yang2024holodeck}
\CSLLeftMargin{{[}21{]} }%
\CSLRightInline{{Y. Yang \emph{et al.}}, {``Holodeck: Language guided
generation of {3D} embodied {AI} environments,''} in \emph{Proceedings
of the IEEE/CVF conference on computer vision and pattern recognition
(CVPR)}, 2024.}

\bibitem[\citeproctext]{ref-gemini31pro}
\CSLLeftMargin{{[}22{]} }%
\CSLRightInline{Google, {``{Gemini 3.1 Pro} (gemini-3.1-pro-preview)
through the gemini {API}.''} 2026. Available:
\url{https://ai.google.dev/}}

\bibitem[\citeproctext]{ref-nanobananapro}
\CSLLeftMargin{{[}23{]} }%
\CSLRightInline{Google, {``{Gemini 3 Pro Image ({`Nano Banana Pro'})
through the Gemini {API}}.''} 2026. Available:
\url{https://ai.google.dev/}}

\bibitem[\citeproctext]{ref-carion2025sam3}
\CSLLeftMargin{{[}24{]} }%
\CSLRightInline{{N. Carion \emph{et al.}}, {``{SAM 3}: {Segment
Anything} with concepts,''} \emph{arXiv preprint arXiv:2511.16719},
2025.}

\bibitem[\citeproctext]{ref-claudecode}
\CSLLeftMargin{{[}25{]} }%
\CSLRightInline{Anthropic, {``{Claude Code}.''} 2026. Available:
\url{https://www.anthropic.com/claude-code}}

\bibitem[\citeproctext]{ref-blender}
\CSLLeftMargin{{[}26{]} }%
\CSLRightInline{Blender Foundation, {``{Blender}.''} 2026. Available:
\url{https://www.blender.org/}}

\bibitem[\citeproctext]{ref-raistrick2026procfunc}
\CSLLeftMargin{{[}27{]} }%
\CSLRightInline{A. Raistrick \emph{et al.}, {``{ProcFunc}:
Function-oriented abstractions for procedural {3D} generation in
{Python},''} \emph{arXiv preprint arXiv:2604.26943}, 2026.}

\bibitem[\citeproctext]{ref-wei2022coacd}
\CSLLeftMargin{{[}28{]} }%
\CSLRightInline{X. Wei, M. Liu, Z. Ling, and H. Su, {``Approximate
convex decomposition for {3D} meshes with collision-aware concavity and
tree search,''} \emph{ACM Transactions on Graphics}, vol. 41, no. 4,
2022.}

\bibitem[\citeproctext]{ref-rempe2026kimodo}
\CSLLeftMargin{{[}29{]} }%
\CSLRightInline{{D. Rempe, M. Petrovich, Y. Yuan, H. Zhang, X. B. Peng,
\emph{et al.}}, {``{Kimodo}: Scaling controllable human motion
generation,''} \emph{arXiv preprint arXiv:2603.15546}, 2026.}

\bibitem[\citeproctext]{ref-motor}
\CSLLeftMargin{{[}30{]} }%
\CSLRightInline{J. P. Kardolus, {``{MOTOR}: Musculoskeletal optimisation
through open robotics simulation.''}
\url{https://github.com/Babon-Innovations-b-v/motor}, 2026.}

\bibitem[\citeproctext]{ref-saito2026soma}
\CSLLeftMargin{{[}31{]} }%
\CSLRightInline{{J. Saito, J. Li, M. de Ruyter, M. Guerrero, \emph{et
al.}}, {``{SOMA}: Unifying parametric human body models,''} \emph{arXiv
preprint arXiv:2603.16858}, 2026.}

\bibitem[\citeproctext]{ref-hicks2015verification}
\CSLLeftMargin{{[}32{]} }%
\CSLRightInline{J. L. Hicks, T. K. Uchida, A. Seth, A. Rajagopal, and S.
L. Delp, {``Is my model good enough? Best practices for verification and
validation of musculoskeletal models and simulations of movement,''}
\emph{Journal of Biomechanical Engineering}, vol. 137, no. 2, p. 020905,
2015.}

\bibitem[\citeproctext]{ref-motorcontact}
\CSLLeftMargin{{[}33{]} }%
\CSLRightInline{J. P. Kardolus, {``{MOTOR} issue 1: Foot contact shares
the load in double support (commit 6d72d73).''}
\url{https://github.com/Babon-Innovations-b-v/motor/issues/1\#issuecomment-6980387981},
2026.}

\end{CSLReferences}

\end{document}
